% ============================================================
%  MF Analytics — Quantitative White Paper
%  Indian Equity Mutual Fund Analytics Framework
% ============================================================
\documentclass[12pt, a4paper]{article}

% ── Packages ────────────────────────────────────────────────
\usepackage[T1]{fontenc}
\usepackage[utf8]{inputenc}
\usepackage{lmodern}
\usepackage{newunicodechar}
\newunicodechar{₹}{\text{\textup{Rs.}}}
\usepackage[margin=2.5cm]{geometry}
\usepackage{amsmath, amssymb, amsthm}
\usepackage{booktabs}
\usepackage{longtable}
\usepackage{array}
\usepackage{multirow}
\usepackage{graphicx}
\usepackage{xcolor}
\usepackage{hyperref}
\usepackage{enumitem}
\usepackage{float}
\usepackage{caption}
\usepackage{subcaption}
\usepackage{mathtools}
\usepackage{bm}
\usepackage{microtype}
\usepackage{parskip}
\usepackage{titlesec}
\usepackage{fancyhdr}
\usepackage{tocloft}
\usepackage{mdframed}
\usepackage{listings}
\usepackage{pdflscape}

% ── Listings style ────────────────────────────────────────────
\lstdefinestyle{pycode}{
  language=Python,
  basicstyle=\ttfamily\footnotesize,
  keywordstyle=\bfseries\color{accent},
  stringstyle=\color{positive},
  commentstyle=\itshape\color{darkgray},
  frame=single,
  backgroundcolor=\color{lightgray},
  breaklines=true,
  breakatwhitespace=true,
  showstringspaces=false,
  tabsize=4,
  captionpos=b,
}
\lstdefinestyle{sqlcode}{
  language=SQL,
  basicstyle=\ttfamily\footnotesize,
  keywordstyle=\bfseries\color{accent},
  frame=single,
  backgroundcolor=\color{lightgray},
  breaklines=true,
  showstringspaces=false,
  tabsize=4,
  captionpos=b,
}

% ── Colours ─────────────────────────────────────────────────
\definecolor{accent}{RGB}{0, 90, 160}
\definecolor{lightgray}{RGB}{245, 245, 245}
\definecolor{darkgray}{RGB}{80, 80, 80}
\definecolor{positive}{RGB}{0, 120, 60}
\definecolor{negative}{RGB}{180, 30, 30}

% ── Hyperref setup ──────────────────────────────────────────
\hypersetup{
  colorlinks = true,
  linkcolor  = accent,
  citecolor  = accent,
  urlcolor   = accent,
  pdftitle   = {MF Analytics: A Quantitative Framework for Indian Equity Mutual Funds},
  pdfauthor  = {MF Analytics},
}

% ── Header / Footer ─────────────────────────────────────────
\pagestyle{fancy}
\fancyhf{}
\fancyhead[L]{\small\textcolor{darkgray}{MF Analytics — Quantitative White Paper}}
\fancyhead[R]{\small\textcolor{darkgray}{\today}}
\fancyfoot[C]{\small\thepage}
\renewcommand{\headrulewidth}{0.4pt}

% ── Section titles ───────────────────────────────────────────
\titleformat{\section}{\large\bfseries\color{accent}}{}{0em}{}[\titlerule]
\titleformat{\subsection}{\normalsize\bfseries\color{accent}}{}{0em}{}
\titleformat{\subsubsection}{\normalsize\bfseries}{}{0em}{}

% ── Theorem-like environments ────────────────────────────────
\newmdtheoremenv[
  backgroundcolor=lightgray,
  linecolor=accent,
  linewidth=1.5pt,
  topline=false, bottomline=false, rightline=false,
]{definition}{Definition}[section]

\newmdtheoremenv[
  backgroundcolor=lightgray,
  linecolor=positive,
  linewidth=1.5pt,
  topline=false, bottomline=false, rightline=false,
]{interpretation}{Interpretation}[section]

% ── Useful macros ────────────────────────────────────────────
\newcommand{\NAV}{\mathrm{NAV}}
\newcommand{\CAGR}{\mathrm{CAGR}}
\newcommand{\rf}{r_f}
\newcommand{\E}{\mathbb{E}}
\newcommand{\Var}{\mathrm{Var}}
\newcommand{\Cov}{\mathrm{Cov}}
\newcommand{\Std}{\mathrm{Std}}
\newcommand{\sign}{\mathrm{sign}}

% ============================================================
\begin{document}

% ── Title Page ──────────────────────────────────────────────
\begin{titlepage}
  \centering
  \vspace*{2cm}
  {\Huge\bfseries\color{accent}
    MF Analytics\\[0.4em]
    \large A Quantitative Framework for\\
    Indian Equity Mutual Fund Analysis
  }\\[2em]
  \rule{0.6\textwidth}{1.2pt}\\[1.5em]
  {\large\textcolor{darkgray}{Quantitative White Paper}}\\[0.5em]
  {\normalsize\textcolor{darkgray}{\today}}\\[3em]

  \begin{abstract}
    \noindent
    This white paper presents the complete quantitative framework underlying the
    MF Analytics platform for Indian equity mutual funds.  We document the
    mathematical definitions, computational procedures, and interpretive
    guidelines for every metric exposed in the platform — spanning return
    analysis, risk measurement, risk-adjusted performance, benchmark-relative
    attribution, Systematic Investment Plan (SIP) modelling, and peer
    comparison via quartile ranking.  All formulas are implemented verbatim in
    the platform's \texttt{analytics.py} engine, making this document both a
    technical reference and an investor education resource.  Where appropriate
    we discuss the economic intuition behind each metric and note common
    misinterpretations.
  \end{abstract}

  \vfill
  {\small\textcolor{darkgray}{
    Data sourced from AMFI India (NAVAll.txt) and Yahoo Finance (index
    prices).  Risk-free rate: India 91-day T-bill, 6.5\% p.a.
    Trading days per year: 252.
  }}
\end{titlepage}

% ── Table of Contents ────────────────────────────────────────
\tableofcontents
\newpage

% ============================================================
\section{Introduction}
% ============================================================

Analysing Indian equity mutual funds rigorously requires navigating a
landscape that is simultaneously rich in data and sparse in standardised
tooling.  The Association of Mutual Funds in India (AMFI) publishes daily
Net Asset Values for over 2{,}000 equity schemes, yet retail investors and
advisors alike typically rely on trailing return tables alone — an
informational diet that obscures risk, benchmark sensitivity, and the
experience of the median systematic investor.

The MF Analytics platform addresses this gap by computing a comprehensive
set of quantitative metrics drawn from modern portfolio theory, empirical
finance, and practitioner risk management.  This white paper formalises
the entire metric library with mathematical precision, documents design
choices made during implementation, and provides interpretive context so
that outputs can be translated into actionable investment insights.

\subsection{Notation Conventions}

Throughout this document we use the following notation.

\begin{center}
\begin{tabular}{@{}ll@{}}
\toprule
Symbol & Meaning \\
\midrule
$P_t$ & NAV (or index level) on trading day $t$ \\
$r_t = P_t/P_{t-1} - 1$ & Simple daily return on day $t$ \\
$\{r_t\}_{t=1}^{T}$ & Time series of $T$ daily returns \\
$\rf$ & Annual risk-free rate (6.5\% p.a.) \\
$r_f^d$ & Daily equivalent: $(1+\rf)^{1/252} - 1$ \\
$N = 252$ & Trading days per year \\
$\mu$ & Mean daily return \\
$\sigma$ & Standard deviation of daily returns \\
$\beta$ & Systematic risk relative to benchmark \\
$R_p, R_b$ & Annualised fund and benchmark returns \\
\bottomrule
\end{tabular}
\end{center}

% ============================================================
\section{Data Infrastructure}
% ============================================================

\subsection{NAV Data Collection}

The platform fetches the AMFI \texttt{NAVAll.txt} file to obtain the
scheme catalogue (scheme code, name, SEBI category, ISIN).  Historical
daily NAVs are then retrieved from \texttt{api.mfapi.in} for each scheme
and stored in a local SQLite database.  Only \emph{Direct Growth} plans
are retained to ensure comparability: regular plans carry distributor
commissions that inflate expense ratios and distort performance
comparisons.

\subsection{Benchmark Index Data}

Index price series are fetched from Yahoo Finance using the
\texttt{yfinance} library.  The platform supports 16 confirmed-working
tickers:

\begin{center}
\begin{tabular}{@{}lll@{}}
\toprule
Index & Ticker & Segment \\
\midrule
Nifty 50 & \texttt{\^{}NSEI} & Broad market large-cap \\
BSE Sensex & \texttt{\^{}BSESN} & Broad market large-cap \\
Nifty Next 50 & \texttt{\^{}NSEMDCP50} & Large/mid transition \\
Nifty Midcap 50 & \texttt{\^{}NSMIDCP} & Mid-cap \\
Nifty Midcap 150 & \texttt{NIFTYMIDCAP150.NS} & Mid-cap \\
Nifty Bank & \texttt{\^{}NSEBANK} & Banking \\
Nifty Financial Svc & \texttt{NIFTY\_FIN\_SERVICE.NS} & Financial services \\
Nifty IT & \texttt{\^{}CNXIT} & Technology \\
Nifty FMCG & \texttt{\^{}CNXFMCG} & Consumer staples \\
Nifty Pharma & \texttt{\^{}CNXPHARMA} & Healthcare \\
Nifty Infrastructure & \texttt{\^{}CNXINFRA} & Infrastructure \\
Nifty Energy & \texttt{\^{}CNXENERGY} & Energy \\
Nifty PSE & \texttt{\^{}CNXPSE} & Public sector enterprises \\
Nifty Auto & \texttt{\^{}CNXAUTO} & Automobiles \& logistics \\
Nifty Realty & \texttt{\^{}CNXREALTY} & Real estate \\
Nifty MNC & \texttt{\^{}CNXMNC} & Multinational companies \\
\bottomrule
\end{tabular}
\end{center}

\subsection{Risk-Free Rate}

All risk-adjusted metrics use a single annualised risk-free rate
$\rf = 6.5\%$, approximating the average India 91-day Treasury Bill yield.
The daily equivalent is computed as
\[
  r_f^d = (1 + \rf)^{1/252} - 1 \approx 0.025\%.
\]

% ============================================================
\section{Fund Classification and Tagging}
% ============================================================

\subsection{SEBI Categories}

SEBI mandates that every open-ended equity mutual fund fall into exactly
one regulatory category.  The platform maps the following categories
directly to a \emph{style tag}:

\begin{center}
\begin{tabular}{@{}ll@{}}
\toprule
SEBI Category & Style Tag \\
\midrule
Equity Scheme -- Large Cap Fund & Large Cap \\
Equity Scheme -- Mid Cap Fund & Mid Cap \\
Equity Scheme -- Small Cap Fund & Small Cap \\
Equity Scheme -- Large \& Mid Cap Fund & Large \& Mid Cap \\
Equity Scheme -- Multi Cap Fund & Multi Cap \\
Equity Scheme -- Flexi Cap Fund & Flexi Cap \\
Equity Scheme -- Focused Fund & Focused \\
Equity Scheme -- Value Fund & Value \\
Equity Scheme -- Contra Fund & Contra \\
Equity Scheme -- ELSS & ELSS \\
Equity Scheme -- Sectoral/Thematic & (resolved by name, see below) \\
\bottomrule
\end{tabular}
\end{center}

\subsection{Thematic Tagging via Rule Engine}

Sectoral and thematic funds do not carry a standardised sub-category
label.  The platform resolves their tag by matching the lowercase fund
name against an ordered list of keyword rules — the first match wins.
Critical ordering constraints are:

\begin{enumerate}[label=\arabic*.]
  \item \textbf{Business Cycle} and \textbf{Momentum} rules precede the
    \textbf{Quant \& Factor} rule, because the AMC ``Quant'' names all
    its funds with the prefix ``quant''; without priority ordering, the
    broad Quant keyword would absorb \emph{quant Business Cycle Fund} and
    \emph{quant Momentum Fund}.
  \item Banking \& Financial Services uses ``banking'', ``financial
    services'', ``bfsi'' but deliberately excludes ``finserv'' to avoid
    false matches on the BAJAJ FINSERV AMC name.
\end{enumerate}

The 26 thematic tags, together with their recommended comparison benchmark,
are listed in the configuration module \texttt{config.TAG\_BENCHMARK}.

% ============================================================
\section{Return Metrics}
% ============================================================

\subsection{Compound Annual Growth Rate (CAGR)}

\begin{definition}[CAGR]
Let $P_0$ and $P_T$ denote the NAV on the first and last day of the
analysis period, and let $D$ be the calendar-day span between them.
\[
  \CAGR = \left(\frac{P_T}{P_0}\right)^{365.25 / D} - 1
\]
\end{definition}

The exponent uses 365.25 to account for leap years.  CAGR is undefined
when $D = 0$ or when $P_T / P_0 \leq 0$ (which cannot occur for NAVs).

\begin{interpretation}[CAGR]
CAGR gives the hypothetical constant annual growth rate that would turn
$P_0$ into $P_T$ in the same time span.  It is a \emph{geometric} average
and therefore correctly reflects compounding.  A fund with CAGR $= 15\%$
doubles in value approximately every 4.8 years (by the rule of 72).
\end{interpretation}

\subsection{Absolute Return}

\[
  R_{\text{abs}} = \frac{P_T}{P_0} - 1
\]

Absolute return is period-length agnostic and should \emph{not} be compared
across funds with different holding periods.  It is shown alongside CAGR
for completeness.

\subsection{Trailing Returns}

For each standard lookback horizon $\Delta$ (1M, 3M, 6M, 1Y, 3Y, 5Y,
10Y), the platform locates the NAV on the date
$t^* = T_{\text{end}} - \Delta$ and computes:

\[
  R_\Delta =
  \begin{cases}
    \dfrac{P_{T_{\text{end}}}}{P_{t^*}} - 1
    & \text{if } \Delta < 1\text{ year (simple return)} \\[1em]
    \left(\dfrac{P_{T_{\text{end}}}}{P_{t^*}}\right)^{365.25/\Delta_d} - 1
    & \text{if } \Delta \geq 1\text{ year (annualised)}
  \end{cases}
\]

where $\Delta_d$ is the lookback in calendar days.  If fewer than 2 NAV
data points exist in the lookback window, the value is reported as missing.

\subsection{Rolling Returns}

For a rolling window of $W$ trading days, the annualised rolling return
at time $t$ is:
\[
  \hat{R}_t^{(W)} = \left(\frac{P_t}{P_{t-W}}\right)^{N/W} - 1
\]

where $N = 252$.  Rolling returns are used to visualise how a fund's return
profile evolves over time and to construct the SIP return distribution.

\subsection{Calendar Year Returns}

The fund's NAV is resampled to year-end values.  For year $y$:
\[
  R_y = \frac{P_{\text{Dec-end}(y)}}{P_{\text{Dec-end}(y-1)}} - 1
\]
The first (and potentially last) calendar year is a partial year; returns
are reported as-is without annualisation.

% ============================================================
\section{Risk Metrics}
% ============================================================

\subsection{Annualised Volatility}

\begin{definition}[Annualised Volatility]
\[
  \sigma = \hat{\sigma}_d \cdot \sqrt{N}
\]
where $\hat{\sigma}_d = \sqrt{\dfrac{1}{T-1}\sum_{t=1}^{T}(r_t - \bar{r})^2}$
is the sample standard deviation of daily returns and $N = 252$.
\end{definition}

\begin{interpretation}[Volatility]
Volatility measures the dispersion of returns around their mean.  A fund
with $\sigma = 20\%$ experiences daily return movements with a typical
magnitude of $20\%/\sqrt{252} \approx 1.26\%$.  Crucially, volatility is
\emph{symmetric} — it penalises upside and downside fluctuations equally.
This is a limitation addressed by the Sortino ratio and the Ulcer Index.
\end{interpretation}

\subsection{Downside Deviation}

\begin{definition}[Downside Deviation]
Let the daily minimum acceptable return be
$\text{MAR}_d = (1 + \rf)^{1/252} - 1$.  Define
$\delta_t = \min(r_t - \text{MAR}_d,\; 0)$.  Then:
\[
  \sigma_D = \sqrt{\frac{1}{T}\sum_{t=1}^{T} \delta_t^2} \cdot \sqrt{N}
\]
\end{definition}

Unlike standard deviation, downside deviation only penalises returns
\emph{below} the minimum acceptable threshold.  It is the denominator
of the Sortino ratio.

\subsection{Maximum Drawdown}

\begin{definition}[Maximum Drawdown]
Let $M_t = \max_{s \leq t} P_s$ be the running all-time-high (ATH) at
time $t$.  The drawdown series is:
\[
  d_t = \frac{P_t - M_t}{M_t} \leq 0
\]
Maximum drawdown is:
\[
  \mathrm{MDD} = \min_t\, d_t
\]
It is reported as a negative number (e.g.\ $-0.35$ means a 35\% peak-to-trough
decline).
\end{definition}

\begin{interpretation}[Maximum Drawdown]
MDD answers: ``what is the worst loss an investor could have experienced if
they bought at the worst possible time?''  For an equity fund MDD below
$-30\%$ is common during bear markets; funds with MDD better than $-20\%$
in such periods often exhibit stronger risk management characteristics.
\end{interpretation}

\subsection{Drawdown Analysis}

Beyond the single worst event, the platform decomposes the entire
drawdown history into \emph{distinct episodes}:

\begin{enumerate}[label=\arabic*.]
  \item An episode begins the first trading day $P_t < M_t$.
  \item The \textbf{trough date} is the day of the minimum $P_t$.
  \item The episode ends when the fund recovers to $P_t \geq M_{\text{start}-1}$
    (a new ATH).  If no recovery has occurred by the data end-date, the
    episode is marked ``ongoing''.
  \item Three durations are recorded:
  \begin{itemize}
    \item \textbf{Total duration} (days from start to recovery)
    \item \textbf{Decline phase} (days from start to trough)
    \item \textbf{Recovery phase} (days from trough to recovery)
  \end{itemize}
\end{enumerate}

Episodes are sorted by depth (worst first).

\subsection{Ulcer Index}

\begin{definition}[Ulcer Index]
\[
  \mathrm{UI} = \sqrt{\frac{1}{T}\sum_{t=1}^{T} d_t^2}
\]
where $d_t$ is the drawdown series defined above.
\end{definition}

The Ulcer Index (Martin \& McCann, 1989) is the root-mean-square of all
drawdown values.  Because $d_t^2$ is large both when drawdowns are deep
\emph{and} when they persist over many days, UI penalises prolonged
underwater periods more severely than MDD does.

\begin{interpretation}[Ulcer Index]
A UI of 5\% is a rough baseline for a well-managed large-cap fund over a
full market cycle.  A fund with UI of 15\% is delivering considerably
more ``ulcer-inducing'' pain to its investors than MDD alone might suggest.
Compare: two funds might share an MDD of $-30\%$, but if one recovered in
60 days while the other spent two years underwater, their UIs will differ
substantially.
\end{interpretation}

\subsection{Average Drawdown}

\[
  \overline{d} = \frac{1}{|\{t: d_t < 0\}|} \sum_{t:\, d_t < 0} d_t
\]

This is the mean drawdown over all days when the fund was below its ATH,
reported as a negative fraction.  It complements MDD by capturing the
typical depth of drawdowns rather than only the worst.

\subsection{Percentage of Time Underwater}

\[
  \text{PTU} = \frac{|\{t : d_t < 0\}|}{T}
\]

The fraction of all trading days on which the fund was below its previous
all-time high.  A PTU of 40\% means the fund was in drawdown 40\% of the
time — common for volatile sector funds.

\subsection{Value at Risk (VaR)}

\begin{definition}[Historical VaR]
At confidence level $\alpha$ (typically 95\% or 99\%), the one-day
Historical VaR is the $\alpha$-quantile of the empirical daily return
distribution:
\[
  \mathrm{VaR}_\alpha = -F^{-1}(1 - \alpha)
\]
where $F$ is the empirical CDF of $\{r_t\}$.  Because VaR is a loss
quantity it is reported as a negative number in the platform (consistent
with the drawdown convention).
\end{definition}

\begin{definition}[Parametric VaR]
Assuming daily returns are normally distributed with mean $\mu$ and
standard deviation $\sigma$:
\[
  \mathrm{VaR}_\alpha^{\text{param}} = \mu + z_{1-\alpha}\,\sigma
\]
where $z_{1-\alpha} = \Phi^{-1}(1-\alpha)$ is the standard normal quantile
(e.g.\ $z_{0.05} = -1.645$ for $\alpha = 0.95$).
\end{definition}

The platform defaults to the historical method, which makes no distributional
assumption.  At least 30 daily observations are required for a meaningful
estimate.

\begin{interpretation}[VaR]
$\mathrm{VaR}_{95} = -2\%$ means: on 95 out of 100 trading days, the fund's
daily loss will not exceed 2\%.  Equivalently, there is a 5\% chance of
losing more than 2\% in a single day.  VaR does \emph{not} describe the
magnitude of losses in the tail — that is the role of CVaR.
\end{interpretation}

\subsection{Conditional VaR (Expected Shortfall)}

\begin{definition}[CVaR / Expected Shortfall]
\[
  \mathrm{CVaR}_\alpha = \E\!\left[r_t \;\middle|\; r_t \leq \mathrm{VaR}_\alpha\right]
  = \frac{1}{|\mathcal{T}|}\sum_{t \in \mathcal{T}} r_t
\]
where $\mathcal{T} = \{t : r_t \leq \mathrm{VaR}_\alpha\}$ is the set of
tail-loss days.
\end{definition}

CVaR is coherent (satisfies sub-additivity) whereas VaR is not.  It
answers: ``given that we are in the worst $(1-\alpha)$\% of outcomes, what
is the \emph{average} loss?''

\subsection{Return Skewness and Kurtosis}

Skewness quantifies the asymmetry of the return distribution:
\[
  \text{Skew} = \frac{\E[(r_t - \mu)^3]}{\sigma^3}
\]
Negative skewness (left tail heavier) is common in equity funds — large
losses are more frequent than large gains of the same magnitude.

Excess kurtosis quantifies tail heaviness relative to the normal
distribution:
\[
  \text{Kurt} = \frac{\E[(r_t - \mu)^4]}{\sigma^4} - 3
\]
Positive excess kurtosis (leptokurtosis) implies more frequent extreme
outcomes (both gains and losses) than a normal distribution would predict.
Most equity return series exhibit positive excess kurtosis of 2--6.

% ============================================================
\section{Risk-Adjusted Return Metrics}
% ============================================================

\subsection{Sharpe Ratio}

\begin{definition}[Sharpe Ratio]
\[
  S = \frac{R_p - \rf}{\sigma}
\]
where $R_p$ is the fund's CAGR, $\rf = 6.5\%$, and $\sigma$ is the
annualised volatility.
\end{definition}

\begin{interpretation}[Sharpe Ratio]
The Sharpe ratio measures the excess return per unit of total risk.  A
value above 1.0 is considered good; above 2.0 is excellent for an equity
fund over a full market cycle.  Values below 0.5 suggest the fund is not
adequately compensating investors for the volatility they bear.

Key limitation: because $\sigma$ is symmetric, a fund that achieves high
returns with high upside volatility (many large positive days) can have an
inflated Sharpe.  This is addressed by the Sortino and Martin ratios.
\end{interpretation}

\subsection{Sortino Ratio}

\begin{definition}[Sortino Ratio]
\[
  \text{Sortino} = \frac{R_p - \rf}{\sigma_D}
\]
where $\sigma_D$ is the annualised downside deviation.
\end{definition}

\begin{interpretation}[Sortino Ratio]
By replacing $\sigma$ with $\sigma_D$, the Sortino ratio only penalises
\emph{negative} deviations from the minimum acceptable return.  A fund
with strongly right-skewed returns (many moderate positives, rare large
negatives) will have Sortino $\gg$ Sharpe, correctly reflecting the
investor's actual experience.  For most equity funds the Sortino ratio is
approximately 1.5--2$\times$ the Sharpe ratio.
\end{interpretation}

\subsection{Calmar Ratio}

\begin{definition}[Calmar Ratio]
\[
  \text{Calmar} = \frac{R_p}{|\mathrm{MDD}|}
\]
\end{definition}

The Calmar ratio (Young, 1991) relates annualised return to the worst
historical drawdown.  It is popular among hedge fund practitioners because
it uses a concrete worst-case loss metric rather than a statistical
quantity.  A Calmar ratio of 0.5 means a fund earned half its worst
drawdown in annualised return — e.g.\ 10\% CAGR against a 20\% MDD.

\subsection{Martin Ratio (Ulcer Performance Index)}

\begin{definition}[Martin Ratio]
\[
  \text{Martin} = \frac{R_p - \rf}{\mathrm{UI}}
\]
where $\mathrm{UI}$ is the Ulcer Index.
\end{definition}

\begin{interpretation}[Martin Ratio]
Proposed by Peter Martin in 1987 (and commercially popularised as the
\emph{Ulcer Performance Index}), the Martin ratio is a Sharpe-like ratio
that substitutes the Ulcer Index for volatility.  Because UI is sensitive
to both depth and duration of drawdowns, the Martin ratio penalises funds
that spend long periods recovering from losses — even if those losses were
individually modest.

Compared to Sharpe: a fund that declines slowly but persistently scores
worse on Martin than Sharpe.  A fund that recovers rapidly from a sharp
crash scores better on Martin than Sharpe.
\end{interpretation}

\subsection{Percentage of Positive Rolling Periods}

\[
  \text{P}^+_{(W)} = \frac{|\{t : \hat{R}_t^{(W)} > 0\}|}{|\{t : \hat{R}_t^{(W)} \text{ defined}\}|}
\]

For a 1-year rolling window ($W = 252$), this answers: ``what fraction of
all possible 1-year holding periods delivered a positive return?''  A
value above 90\% indicates the fund rarely disappoints investors who hold
for a full year.

% ============================================================
\section{Benchmark-Relative Metrics}
% ============================================================

Let $\{r_t^f\}$ and $\{r_t^b\}$ denote the daily returns of the fund and
benchmark, respectively, on their common trading days.  All benchmark
metrics require at least 30 aligned observations.

\subsection{Beta}

\begin{definition}[Beta]
\[
  \beta = \frac{\Cov(r^f, r^b)}{\Var(r^b)}
\]
\end{definition}

Beta measures the fund's systematic sensitivity to benchmark movements.
A beta of 1.2 implies the fund tends to move 1.2\% for every 1\% benchmark
move (leveraged market exposure).  Beta below 1 implies defensive
characteristics.

The sample estimator used is:
\[
  \hat{\beta} = \frac{\sum_{t=1}^{T}(r_t^f - \bar{r}^f)(r_t^b - \bar{r}^b)}
               {\sum_{t=1}^{T}(r_t^b - \bar{r}^b)^2}
\]

\subsection{Jensen's Alpha}

\begin{definition}[Jensen's Alpha]
\[
  \alpha = R_p - \left[\rf + \beta\,(R_b - \rf)\right]
\]
where $R_p$ and $R_b$ are the fund's and benchmark's annualised CAGR.
\end{definition}

\begin{interpretation}[Jensen's Alpha]
Alpha is the return earned in excess of what the Capital Asset Pricing
Model (CAPM) predicts given the fund's market exposure.  Positive alpha
indicates skill: the fund manager added value beyond simply taking
benchmark risk.

Numerically, $\alpha = 3\%$ means the fund delivered 3 percentage points
per year more than a passive strategy with the same beta.  The platform
computes alpha using the full-period CAGR for both the fund and benchmark,
rather than OLS regression intercepts, to keep the interpretation consistent
with the other return metrics.
\end{interpretation}

\subsection{Tracking Error}

\begin{definition}[Tracking Error]
\[
  \mathrm{TE} = \Std(r_t^f - r_t^b) \cdot \sqrt{N}
\]
\end{definition}

Tracking error is the annualised standard deviation of daily active returns
(fund minus benchmark).  For index funds, TE should be below 0.5\%; for
actively managed funds TE typically ranges from 3\% to 10\%.

\subsection{Information Ratio}

\begin{definition}[Information Ratio]
\[
  \mathrm{IR} = \frac{R_p - R_b}{\mathrm{TE}}
\]
\end{definition}

The Information Ratio (IR) measures the consistency of outperformance: it
rewards a high active return ($R_p - R_b$) but penalises inconsistency
(high TE).  An IR above 0.5 over a 3-year period is widely considered
evidence of skill.  The IR is often called the ``Sharpe ratio of active
management''.

\subsection{Treynor Ratio}

\begin{definition}[Treynor Ratio]
\[
  T = \frac{R_p - \rf}{\beta}
\]
\end{definition}

The Treynor ratio replaces total volatility with systematic risk ($\beta$).
It is most useful when comparing funds that are part of a well-diversified
portfolio, where idiosyncratic risk is already diversified away.

\subsection{Upside and Downside Capture Ratios}

Capture ratios are computed on \emph{monthly} returns.  Let
$\mathcal{U} = \{m : r_m^b > 0\}$ be the set of months in which the
benchmark was positive (``up months''), and
$\mathcal{D} = \{m : r_m^b < 0\}$ the set of ``down months''.

\begin{definition}[Upside Capture Ratio]
\[
  \mathrm{UCR} = \frac{G_{\mathcal{U}}^f}{G_{\mathcal{U}}^b} \times 100
\]
where $G_{\mathcal{U}}^f = \left(\prod_{m \in \mathcal{U}}(1 + r_m^f)\right)^{1/|\mathcal{U}|} - 1$
is the geometric mean monthly return of the fund in up months, and
similarly for $G_{\mathcal{U}}^b$.
\end{definition}

\begin{definition}[Downside Capture Ratio]
\[
  \mathrm{DCR} = \frac{G_{\mathcal{D}}^f}{G_{\mathcal{D}}^b} \times 100
\]
\end{definition}

\begin{interpretation}[Capture Ratios]
An ideal fund has UCR $> 100$ (participates fully in rallies) and
DCR $< 100$ (cushions declines).  A fund with UCR $= 110$ and DCR $= 80$
captures 110\% of up-market gains but only 80\% of down-market losses —
an asymmetry highly favourable to long-term compounding.  Funds in the
upper-left quadrant of a UCR vs DCR scatter plot are the most desirable.

Note: capture ratios below 3 months of data are unreliable and
marked as missing.
\end{interpretation}

\subsection{Batting Average}

\begin{definition}[Batting Average]
\[
  \mathrm{BA} = \frac{|\{m : r_m^f > r_m^b\}|}{M}
\]
where $M$ is the total number of months with data for both the fund and
benchmark.
\end{definition}

Batting average measures the frequency of outperformance, without regard
to magnitude.  A BA of 0.55 means the fund beat its benchmark in 55\% of
months.  A high BA combined with a high UCR and low DCR is a strong
signal of consistent active management quality.

% ============================================================
\section{Rolling Metrics}
% ============================================================

All rolling metrics use a backward-looking window of $W$ trading days
(configurable: 126, 252, 504, 756, 1260 for 6M, 1Y, 2Y, 3Y, 5Y).

\subsection{Rolling Volatility}

\[
  \sigma_t^{(W)} = \hat{\sigma}\!\left(\{r_{t-W+1}, \ldots, r_t\}\right) \cdot \sqrt{N}
\]

\subsection{Rolling Sharpe Ratio}

\[
  S_t^{(W)} = \frac{\hat{R}_t^{(W)} - \rf}{\sigma_t^{(W)}}
\]

where $\hat{R}_t^{(W)} = (P_t/P_{t-W})^{N/W} - 1$ is the annualised
return over the trailing window.

\subsection{Rolling Sortino Ratio}

\[
  \text{Sortino}_t^{(W)} =
  \frac{\hat{R}_t^{(W)} - \rf}{\hat{\sigma}_{D,t}^{(W)}}
\]

where $\hat{\sigma}_{D,t}^{(W)}$ is the rolling downside deviation:
\[
  \hat{\sigma}_{D,t}^{(W)} = \sqrt{\frac{1}{W}\sum_{s=t-W+1}^{t}\min(r_s - r_f^d, 0)^2} \cdot \sqrt{N}
\]

\subsection{Rolling Alpha and Beta}

\[
  \hat{\beta}_t^{(W)} = \frac{\sum_{s=t-W+1}^{t}(r_s^f - \bar{r}^f)(r_s^b - \bar{r}^b)}
                          {\sum_{s=t-W+1}^{t}(r_s^b - \bar{r}^b)^2}
\]

\[
  \hat{\alpha}_t^{(W)} = \hat{R}_t^{(W),f} - \left[\rf + \hat{\beta}_t^{(W)}\left(\hat{R}_t^{(W),b} - \rf\right)\right]
\]

Rolling alpha and beta reveal structural changes in a manager's risk
positioning over time.  A declining beta over time may signal a defensive
tilt, while rising alpha may indicate an improving stock-selection process.

% ============================================================
\section{SIP Analysis and XIRR}
% ============================================================

Systematic Investment Plans (SIPs) are the dominant mode of retail mutual
fund investing in India.  Evaluating a fund purely on point-to-point CAGR
ignores the cost averaging and path-dependency inherent in monthly
investments.

\subsection{XIRR: Extended Internal Rate of Return}

XIRR is the standard of AMFI, SEBI disclosures, and Excel for evaluating
irregular cashflow series.  It solves for the annualised discount rate $r$
that sets the Net Present Value of all cashflows to zero.

\begin{definition}[XIRR]
Given cashflows $C_0, C_1, \ldots, C_n$ at dates $t_0 < t_1 < \cdots < t_n$
(with $t_0$ as the reference date), XIRR solves:
\[
  \sum_{i=0}^{n} \frac{C_i}{(1 + r)^{(t_i - t_0)/365.25}} = 0
\]
Sign convention (Excel XIRR): investments are \emph{negative} (cash
out of the investor's pocket); redemptions are \emph{positive} (cash
received).
\end{definition}

The platform uses \texttt{scipy.optimize.brentq} with bracket
$r \in (-0.9999,\; 200.0)$ to accommodate the very high XIRRs achievable
by short-term SIPs in bull markets.

\subsection{SIP Return Computation}

For a $Y$-year SIP ending on the last available NAV date:

\begin{enumerate}[label=\arabic*.]
  \item Generate $Y \times 12$ month-start dates ending at the last NAV date.
  \item For each month-start date $m_i$, find the NAV on or immediately after
    $m_i$ in the price series.
  \item Invest ₹1 at that NAV, purchasing $u_i = 1/\NAV_i$ units:
    \[
      C_i = -1, \quad \text{units accumulated} \;\to\; U = \sum_i u_i
    \]
  \item At the end of the period, redeem all units at the final NAV:
    \[
      C_{\text{redemption}} = +U \times P_T
    \]
  \item Solve $\sum_i C_i/(1+r)^{(t_i-t_0)/365.25} = 0$ for $r$.
\end{enumerate}

SIP XIRR is computed for horizons $Y \in \{1, 3, 5, 10\}$ years.

\begin{interpretation}[SIP XIRR]
SIP XIRR captures cost averaging: if a fund fell sharply midway through
the SIP and then recovered, more units were accumulated at low prices,
and the XIRR can exceed the equivalent lump-sum CAGR.  Conversely, a
fund that rallied early and then corrected will show SIP XIRR below
its CAGR.  For most long-term equity funds, 5-year and 10-year SIP
XIRR is a robust estimate of the ``real investor experience''.
\end{interpretation}

\subsection{Rolling SIP Distribution}

To assess \emph{consistency} of SIP outcomes, the platform evaluates XIRR
for every possible $Y$-year SIP starting date in the fund's history.
This produces a distribution of XIRRs indexed by start date.  Summary
statistics of interest:

\begin{itemize}
  \item \textbf{Percentage positive}: fraction of all $Y$-year SIP windows
    with XIRR $> 0$.  Values below 90\% are concerning for a 3-year
    horizon and should be near 100\% for a 10-year horizon in a mature
    equity fund.
  \item \textbf{Worst XIRR}: the SIP window that started at the most
    unfavourable time (typically right before a major crash).
  \item \textbf{Median XIRR}: more robust than mean to outlier windows.
\end{itemize}

% ============================================================
\section{Peer Comparison and Quartile Ranking}
% ============================================================

\subsection{Metrics Comparison Table}

The platform constructs an $M \times F$ matrix where $M$ metrics are
rows and $F$ funds (within a tag or user-selected group) are columns.
Each cell $(m, f)$ contains $v_{m,f}$, the value of metric $m$ for fund
$f$.  Values are formatted as percentages for return/risk metrics and
as plain decimals for ratios.

\subsection{Within-Category Quartile Ranking}

For each metric row $m$, funds are ranked against their peers:

\begin{enumerate}[label=\arabic*.]
  \item Let $\mathcal{F}_m = \{f : v_{m,f} \text{ is finite}\}$ be the
    set of funds with valid values.
  \item Determine the ranking direction:
    \[
      \text{ascending} =
      \begin{cases}
        \text{True} & \text{if } m \in \mathcal{L} \quad \text{(lower is better)} \\
        \text{False} & \text{otherwise (higher is better)}
      \end{cases}
    \]
    where $\mathcal{L}$ = \{volatility, ulcer\_index, average\_drawdown,
    pct\_time\_underwater, tracking\_error, var\_95, var\_99, cvar\_95,
    cvar\_99, worst\_drawdown\_depth, worst\_drawdown\_duration,
    downside\_capture\}.
  \item Assign rank $\rho_f = \text{rank}(v_{m,f}) / |\mathcal{F}_m|$
    (using minimum rank for ties).
  \item Classify:
    \[
      Q_f =
      \begin{cases}
        \text{Q1} & \rho_f \leq 0.25 \\
        \text{Q2} & 0.25 < \rho_f \leq 0.50 \\
        \text{Q3} & 0.50 < \rho_f \leq 0.75 \\
        \text{Q4} & \rho_f > 0.75
      \end{cases}
    \]
\end{enumerate}

Q1 (best quartile) is rendered in green; Q4 (worst quartile) in red.

\begin{interpretation}[Quartile Rankings]
Quartile rankings provide the answer to the question that raw numbers
cannot: ``is this fund good \emph{relative to its peers}?''  A
large-cap fund with a Sharpe of 0.7 may look mediocre in absolute terms
but could be Q1 within its category if most large-cap funds have Sharpe
below 0.5.  Conversely, a sector fund with Sharpe 1.5 might be Q3 if it
operates in a high-return, low-volatility niche.

Users should look for funds that are consistently Q1 or Q2 across
\emph{multiple} metrics rather than exceptional in one metric and poor
in others.
\end{interpretation}

% ============================================================
\section{Return Colour Coding}
% ============================================================

The heatmap of trailing returns uses a signed per-column gradient:

\[
  \text{scaled}_{f,\Delta} = \frac{R_{f,\Delta}/\max_f |R_{f,\Delta}| + 1}{2}
  \;\in\; [0, 1]
\]

This maps the range $[-\max|\cdot|, 0, +\max|\cdot|]$ to $[0, 0.5, 1]$
in the \texttt{RdYlGn} colormap: red for negative returns, yellow for zero,
and green for positive returns — regardless of whether all values happen to
be positive or all negative.  The per-column normalisation ensures that both
short (1M) and long (5Y) columns use the full colour range, making it easy
to spot anomalies at any horizon.

% ============================================================
\section{Implementation Notes}
% ============================================================

\subsection{Data Quality}

\begin{itemize}
  \item Dates in \texttt{mfapi.in} are in DD-MMM-YYYY format (e.g.\
    ``15-Jan-2020'') and are converted to ISO 8601 before storage.
  \item Duplicate NAV entries for the same (scheme\_code, date) are
    silently deduplicated (last-write wins).
  \item The benchmark series and fund NAV series are aligned on common
    dates before any benchmark-relative computation.
  \item All metrics require a minimum of 30 daily observations; fewer
    than 252 points disables rolling 1-year metrics.
\end{itemize}

\subsection{Numerical Stability}

\begin{itemize}
  \item The XIRR solver uses Brent's method with bracket
    $r \in (-0.9999, 200.0)$ and tolerance $10^{-6}$.  A bracket of 200
    (20{,}000\% p.a.) is required because a 1-year SIP into a 4$\times$
    equity rally can legitimately produce XIRRs above 100\%.
  \item Rolling beta denominators of zero (flat benchmark windows) are
    replaced by NaN.
  \item All downstream callers test for NaN and $\pm\infty$ using
    \texttt{pd.notna(v) and np.isfinite(v)} rather than
    \texttt{np.isnan(v)}, which raises \texttt{TypeError} on
    \texttt{None}.
\end{itemize}

\subsection{Tagger Rule Ordering}

The thematic keyword rule engine uses a first-match-wins strategy.  The
ordering constraint that is non-obvious and critical: \emph{Business Cycle}
and \emph{Momentum} rules must appear before the \emph{Quant \& Factor} rule
in the ordered list.  Quant Asset Management Company names all its funds
``quant XYZ Fund'', so any keyword containing ``quant'' would absorb
``quant Business Cycle Fund'' and ``quant Momentum Fund'' into the wrong
bucket.  The solution is to restrict Quant \& Factor keywords to
``quant fund'', ``quantamental'', etc.\ (not bare ``quant''), and ensure
Business Cycle and Momentum patterns match first.

\subsection{Common-Period Calculation}

When displaying all funds in a tag on a single growth chart, a common
analysis window is required.  Naively using $[\max(\text{starts}),
\min(\text{ends})]$ fails when a tag contains a newly-launched fund
(start $\approx$ today) alongside discontinued funds (end in the past),
yielding an empty or inverted window.

The platform uses:
\[
  t_{\text{end}} = \max(\text{all end dates}), \qquad
  t_{\text{start}} = P_{80}\!\left(\text{sorted start dates}\right)
\]

where $P_{80}$ is the 80th percentile of start dates.  This excludes the
20\% most-recently-launched funds from constraining the window, while still
representing the majority of funds over a meaningful history.

% ============================================================
\section{Summary Metric Reference}
% ============================================================

\begin{longtable}{@{}p{3.8cm}p{5cm}p{5.2cm}@{}}
\caption{Complete Metric Reference}\\
\toprule
\textbf{Metric} & \textbf{Formula / Definition} & \textbf{Interpretation} \\
\midrule
\endfirsthead
\multicolumn{3}{c}{\textit{(continued from previous page)}}\\
\toprule
\textbf{Metric} & \textbf{Formula / Definition} & \textbf{Interpretation} \\
\midrule
\endhead
\midrule
\multicolumn{3}{r}{\textit{continued on next page}}\\
\endfoot
\bottomrule
\endlastfoot

CAGR
  & $(P_T/P_0)^{365.25/D} - 1$
  & Annualised total return; primary headline metric. \\[4pt]

Trailing Returns
  & Period-end vs period-start; annualised for $\geq$1Y
  & Snapshot of recent performance at standard horizons. \\[4pt]

Volatility
  & $\hat\sigma_d \cdot \sqrt{252}$
  & Total risk; symmetric — penalises up and down moves. \\[4pt]

Downside Dev.
  & $\sqrt{\tfrac{1}{T}\sum\min(r_t - r_f^d, 0)^2} \cdot \sqrt{252}$
  & Risk below MAR; denominator of Sortino ratio. \\[4pt]

Max Drawdown
  & $\min_t (P_t - M_t)/M_t$
  & Worst peak-to-trough decline; tail-risk proxy. \\[4pt]

Ulcer Index
  & $\sqrt{\tfrac{1}{T}\sum d_t^2}$
  & RMS of all drawdowns; penalises depth \emph{and} duration. \\[4pt]

Average Drawdown
  & Mean of $d_t$ where $d_t < 0$
  & Typical drawdown depth while underwater. \\[4pt]

\% Time Underwater
  & $|\{t : d_t < 0\}| / T$
  & Fraction of days below all-time high. \\[4pt]

VaR (95\%, 99\%)
  & $\alpha$-quantile of $\{r_t\}$ (historical)
  & Daily loss not exceeded on $(1-\alpha)$ of days. \\[4pt]

CVaR (95\%, 99\%)
  & $\mathbb{E}[r_t | r_t \leq \mathrm{VaR}_\alpha]$
  & Expected loss in the worst $(1-\alpha)$\% of days. \\[4pt]

Skewness
  & $\mathbb{E}[(r_t-\mu)^3] / \sigma^3$
  & Negative = fatter left tail (more frequent large losses). \\[4pt]

Kurtosis
  & $\mathbb{E}[(r_t-\mu)^4] / \sigma^4 - 3$
  & Positive = heavier tails than normal distribution. \\[4pt]

Sharpe Ratio
  & $(R_p - r_f) / \sigma$
  & Excess return per unit of total risk. $>1$ is good. \\[4pt]

Sortino Ratio
  & $(R_p - r_f) / \sigma_D$
  & Excess return per unit of downside risk. \\[4pt]

Calmar Ratio
  & $R_p / |\mathrm{MDD}|$
  & Return vs worst loss; higher is better. \\[4pt]

Martin Ratio
  & $(R_p - r_f) / \mathrm{UI}$
  & Sharpe variant using Ulcer Index; penalises prolonged pain. \\[4pt]

\% Positive 1Y
  & Fraction of 1Y rolling windows with $>0$ return
  & Consistency of annual positive returns. \\[4pt]

Beta
  & $\mathrm{Cov}(r^f, r^b) / \mathrm{Var}(r^b)$
  & Systematic exposure to benchmark; 1 = market-neutral risk. \\[4pt]

Jensen's Alpha
  & $R_p - [r_f + \beta(R_b - r_f)]$
  & Excess return beyond CAPM; positive = manager skill. \\[4pt]

Tracking Error
  & $\mathrm{Std}(r^f - r^b) \cdot \sqrt{252}$
  & Variability of active returns; high for concentrated bets. \\[4pt]

Information Ratio
  & $(R_p - R_b) / \mathrm{TE}$
  & Consistency of outperformance; $>0.5$ signals skill. \\[4pt]

Treynor Ratio
  & $(R_p - r_f) / \beta$
  & Return per unit of \emph{systematic} risk only. \\[4pt]

Upside Capture
  & $(G_\mathcal{U}^f / G_\mathcal{U}^b) \times 100$
  & $>100$: fund rises more than benchmark in up months. \\[4pt]

Downside Capture
  & $(G_\mathcal{D}^f / G_\mathcal{D}^b) \times 100$
  & $<100$: fund falls less than benchmark in down months. \\[4pt]

Batting Average
  & Fraction of months fund $>$ benchmark
  & Frequency of monthly outperformance. \\[4pt]

SIP XIRR 1Y/3Y/5Y/10Y
  & Brent's method on monthly ₹1 SIP cashflows
  & Annualised return of the systematic investor. \\[4pt]

\end{longtable}

% ============================================================
\section{Common Misinterpretations and Caveats}
% ============================================================

\begin{enumerate}[label=\textbf{\arabic*.}]

\item \textbf{Short-period CAGR is misleading.}  A fund with 6 months of
  history showing CAGR 40\% simply experienced a bull market at launch.
  Require at least 3 years of history before drawing conclusions about
  manager skill.

\item \textbf{Alpha without statistical testing.}  The platform reports
  Jensen's alpha as a point estimate.  A positive alpha does not
  necessarily reflect skill — it could be luck, style exposure, or
  small-sample noise.  For robustness, look for consistent positive
  alpha across multiple sub-periods.

\item \textbf{Beta is not constant.}  Fund managers actively change their
  market exposure.  A fund's full-history beta may differ substantially
  from its recent 1-year rolling beta.  Always consult the rolling beta
  chart alongside the static value.

\item \textbf{VaR does not bound the loss.}  VaR$_{95}$ = $-2\%$ means
  losses exceed 2\% in 5\% of days — but on those bad days, losses
  could be $-5\%$, $-10\%$, or worse.  CVaR (Expected Shortfall)
  provides the expected loss in that tail.

\item \textbf{Tracking error and alpha are not independent.}  A manager
  with TE = 0.5\% cannot have large alpha — they are effectively
  running an index fund.  High TE is the \emph{necessary} (but not
  sufficient) condition for large positive alpha.

\item \textbf{SIP XIRR depends on the end date.}  If the fund's current
  NAV is at a cyclical peak, the SIP XIRR will be inflated; if it is in
  a trough, it will be depressed.  The rolling SIP distribution is a
  more robust characterisation of the likely investor experience.

\item \textbf{Quartile ranks are peer-relative.}  A Q1 fund in a poor
  category may still be an absolute underperformer.  Always examine
  absolute values alongside peer ranks.

\item \textbf{Direct vs Regular plan comparison.}  The platform
  exclusively uses Direct Growth plans.  Comparing results against
  industry data that mixes Direct and Regular plans will show an
  apparent outperformance equal approximately to the TER differential
  (typically 0.5--1.5\% p.a.).

\end{enumerate}

% ============================================================
\section{Platform Architecture and Implementation}
\label{sec:architecture}
% ============================================================

This section documents the full technical architecture of the MF Analytics
platform — covering module design, database schema, data pipeline, application
layer, tagging engine, caching strategy, and deployment.  Together with the
metric definitions in earlier sections, it constitutes a complete specification
of the system.

% -----------------------------------------------------------
\subsection{High-Level Architecture}
% -----------------------------------------------------------

The platform is a single-host, Python-based application with four logical
layers:

\begin{center}
\begin{tabular}{@{}p{0.22\textwidth}p{0.70\textwidth}@{}}
\toprule
\textbf{Layer} & \textbf{Components} \\
\midrule
\textbf{Presentation} & Streamlit web application (\texttt{app.py}) —
  multi-page UI rendered in the browser \\[4pt]
\textbf{Analytics} & Pure-Python computation engine (\texttt{analytics.py}) —
  stateless functions operating on \texttt{pd.Series} inputs \\[4pt]
\textbf{Data Access} & SQLite database (\texttt{mf\_data.db}) accessed through
  \texttt{db.py}; fetchers in \texttt{fetcher.py} bridge external APIs to the
  local store \\[4pt]
\textbf{Configuration} & \texttt{config.py} — all environment-specific
  constants (paths, URLs, rates, benchmarks) in one place \\
\bottomrule
\end{tabular}
\end{center}

\begin{mdframed}[backgroundcolor=lightgray,linecolor=accent,linewidth=1.5pt,
                 topline=false,bottomline=false,rightline=false]
\textbf{Design principle.}
The analytics layer is fully decoupled from both the database and the UI:
\texttt{analytics.py} accepts plain \texttt{pd.Series} arguments and returns
plain Python dicts or Series.  This makes every metric independently testable
and reusable outside Streamlit.
\end{mdframed}

\subsubsection{Module dependency graph}

\begin{center}
\begin{tabular}{@{}l@{}}
\toprule
\ttfamily
\begin{tabular}{@{}l@{}}
config.py \\
\quad$\uparrow$ imported by all modules \\[4pt]
db.py \quad (depends on: config) \\
fetcher.py \quad (depends on: config, db) \\
tagger.py \quad (depends on: db) \\
analytics.py \quad (depends on: config only) \\
update.py \quad (depends on: db, fetcher) \\
app.py \quad (depends on: config, db, fetcher, tagger, analytics) \\
\end{tabular}\\
\bottomrule
\end{tabular}
\end{center}

% -----------------------------------------------------------
\subsection{Source Code Modules}
% -----------------------------------------------------------

\begin{longtable}{@{}p{0.20\textwidth}r p{0.60\textwidth}@{}}
\caption{Module inventory}\label{tab:modules}\\
\toprule
\textbf{Module} & \textbf{LOC} & \textbf{Responsibility} \\
\midrule
\endfirsthead
\multicolumn{3}{c}{\textit{(continued)}}\\
\toprule
\textbf{Module} & \textbf{LOC} & \textbf{Responsibility} \\
\midrule
\endhead
\texttt{config.py} & 132 &
  All configuration constants: file paths, API endpoints, rate limits,
  risk-free rate, trading-days constant, benchmark ticker map,
  NSE direct-index map, tag-to-benchmark map, rolling windows,
  VaR confidence levels. \\[4pt]
\texttt{db.py} & 260 &
  SQLite schema definition, WAL-mode connection context manager,
  upsert/bulk-insert helpers, query helpers for NAV and benchmark
  retrieval, update-log management, and DB statistics. \\[4pt]
\texttt{fetcher.py} & 511 &
  Data acquisition from three external sources: AMFI
  \texttt{NAVAll.txt} (scheme list), \texttt{mfapi.in} REST API
  (full NAV history per scheme), \texttt{yfinance} (benchmark index
  prices), and the NSE daily archive CSV files (smallcap/microcap
  indices not available on yfinance). \\[4pt]
\texttt{tagger.py} & 317 &
  Two-pass rule engine that assigns sector/style tags to every fund.
  Manages the \texttt{fund\_tags} table, exposes query helpers for
  tag-based filtering, and supports per-fund manual tag override. \\[4pt]
\texttt{analytics.py} & 765 &
  Complete metric library: trailing returns, CAGR, volatility,
  drawdown series, Sharpe/Sortino/Calmar/Martin/Ulcer ratios,
  alpha, beta, $R^2$, tracking error, information ratio,
  capture ratios, rolling variants of each, SIP XIRR,
  calendar year returns, peer comparison table. \\[4pt]
\texttt{update.py} & 113 &
  Standalone daily updater script.  Orchestrates scheme-list refresh,
  incremental NAV download, and benchmark refresh.  Supports
  one-shot and scheduled (\texttt{cron}-friendly) execution modes. \\[4pt]
\texttt{app.py} & 2056 &
  Streamlit application.  Five pages (Dashboard, Fund Analysis,
  Fund Comparison, Sector \& Style, Data Management), module-level
  \texttt{@st.cache\_data} wrappers, sidebar navigation, and all
  interactive visualisations. \\
\bottomrule
\end{longtable}

% -----------------------------------------------------------
\subsection{Database Schema}
% -----------------------------------------------------------

All persistent state lives in a single SQLite file (\texttt{mf\_data.db})
co-located with the source code.  WAL journal mode and \texttt{PRAGMA
synchronous=NORMAL} are set on every connection for safe concurrent reads while
a write is in progress.

\subsubsection{Table: \texttt{funds}}

\begin{lstlisting}[style=sqlcode,caption={funds — scheme master}]
CREATE TABLE funds (
    scheme_code   INTEGER PRIMARY KEY,   -- AMFI unique scheme code
    scheme_name   TEXT    NOT NULL,
    scheme_category TEXT,               -- e.g. "Equity Scheme - Large Cap Fund"
    scheme_type   TEXT,                 -- "Open Ended Schemes" etc.
    fund_house    TEXT,                 -- derived from scheme_name (first token)
    isin_growth   TEXT,                 -- ISIN of growth variant
    isin_div      TEXT                  -- ISIN of IDCW/dividend variant
);
CREATE INDEX idx_funds_category ON funds(scheme_category);
\end{lstlisting}

The table is populated from AMFI \texttt{NAVAll.txt} on every scheme-list
refresh.  Upserts use \texttt{ON CONFLICT … DO UPDATE SET} with
\texttt{COALESCE} so that non-null fields are never overwritten by a null.

\subsubsection{Table: \texttt{nav\_history}}

\begin{lstlisting}[style=sqlcode,caption={nav\_history — daily NAV time series}]
CREATE TABLE nav_history (
    scheme_code  INTEGER NOT NULL,
    date         TEXT    NOT NULL,   -- ISO-8601: "YYYY-MM-DD"
    nav          REAL    NOT NULL,
    PRIMARY KEY (scheme_code, date),
    FOREIGN KEY (scheme_code) REFERENCES funds(scheme_code)
);
CREATE INDEX idx_nav_date   ON nav_history(date);
CREATE INDEX idx_nav_scheme ON nav_history(scheme_code);
\end{lstlisting}

Records are inserted with \texttt{INSERT OR IGNORE} so incremental downloads
never duplicate data.  Date strings are stored in ISO-8601 lexicographic order
so \texttt{MAX(date)} and range queries work without date parsing.

\subsubsection{Table: \texttt{benchmark\_data}}

\begin{lstlisting}[style=sqlcode,caption={benchmark\_data — index close prices}]
CREATE TABLE benchmark_data (
    index_name   TEXT NOT NULL,   -- display name, e.g. "Nifty 50"
    date         TEXT NOT NULL,   -- ISO-8601: "YYYY-MM-DD"
    close_price  REAL NOT NULL,
    PRIMARY KEY (index_name, date)
);
CREATE INDEX idx_benchmark_date ON benchmark_data(date);
\end{lstlisting}

Both yfinance-sourced indices and NSE-archive-sourced indices are stored in
this single table; the \texttt{index\_name} column holds the display name used
throughout the UI.

\subsubsection{Table: \texttt{fund\_tags}}

\begin{lstlisting}[style=sqlcode,caption={fund\_tags — sector/style taxonomy}]
CREATE TABLE fund_tags (
    scheme_code INTEGER NOT NULL,
    tag         TEXT    NOT NULL,
    PRIMARY KEY (scheme_code, tag),
    FOREIGN KEY (scheme_code) REFERENCES funds(scheme_code)
);
CREATE INDEX idx_fund_tags_tag  ON fund_tags(tag);
CREATE INDEX idx_fund_tags_code ON fund_tags(scheme_code);
\end{lstlisting}

Tags are assigned automatically by the tagger rule engine (see
Section~\ref{sec:tagger}) and can be overridden per-fund via the Data
Management page.

\subsubsection{Table: \texttt{update\_log}}

\begin{lstlisting}[style=sqlcode,caption={update\_log — audit trail}]
CREATE TABLE update_log (
    id              INTEGER PRIMARY KEY AUTOINCREMENT,
    update_type     TEXT    NOT NULL,   -- "nav", "benchmark", etc.
    started_at      TEXT    NOT NULL,
    finished_at     TEXT,
    schemes_updated INTEGER DEFAULT 0,
    status          TEXT    DEFAULT 'running'  -- 'completed' | 'error'
);
\end{lstlisting}

Every update run writes a row on start and updates it on completion.  The
Data Management page reads the latest row to display last-update status and
timestamp.

% -----------------------------------------------------------
\subsection{Data Pipeline}
% -----------------------------------------------------------

Data enters the system in four stages, each handled by a function in
\texttt{fetcher.py}.

\begin{center}
\begin{tabular}{@{}l@{\;$\longrightarrow$\;}l@{\;$\longrightarrow$\;}l@{}}
AMFI NAVAll.txt & mfapi.in REST API & SQLite \texttt{nav\_history} \\
NSE Archive CSV & NSE bulk ZIP downloads & SQLite \texttt{benchmark\_data} \\
Yahoo Finance & yfinance Python library & SQLite \texttt{benchmark\_data} \\
\end{tabular}
\end{center}

\subsubsection{Stage 1 — Scheme Discovery (\texttt{fetch\_scheme\_list})}

AMFI publishes a plain-text file at
\texttt{amfiindia.com/spages/NAVAll.txt} that lists every registered mutual
fund scheme in India.  Each line is semicolon-delimited:

\begin{center}
\texttt{SchemeCode ; ISIN\_Div ; ISIN\_Growth ; SchemeName ; NAV ; Date}
\end{center}

Section headers of the form \texttt{Open Ended Schemes(Equity Scheme - Large
Cap Fund)} identify the SEBI category for the schemes that follow.
The parser (\texttt{parse\_amfi\_nav\_file}) maintains a running
\texttt{current\_category} variable updated on each header line, then attaches
it to every data record.

The filter \texttt{filter\_equity\_schemes} retains only Direct Growth plans
in equity or ELSS categories.  It rejects IDCW, dividend, payout, regular,
institutional, and non-equity scheme names via a keyword exclusion list
(\texttt{\_DIRECT\_GROWTH\_EXCLUDE}).

\subsubsection{Stage 2 — NAV History (\texttt{download\_nav\_for\_schemes})}

Full NAV history for each scheme is fetched from the open \texttt{mfapi.in}
REST API:

\begin{center}
\texttt{GET https://api.mfapi.in/mf/\{scheme\_code\}}
\end{center}

The API returns up to the full history in a single JSON response.  Each record
contains a date string in \texttt{dd-mm-yyyy} format which is converted to
ISO-8601 (\texttt{yyyy-mm-dd}) before storage.  The download is incremental:
\texttt{get\_last\_nav\_date} retrieves the most recent date already stored and
only records newer than that date are inserted.

Rate limiting is enforced at \texttt{MFAPI\_DELAY\_SECONDS} (0.5 s) between
requests and database commits happen every \texttt{MFAPI\_BATCH\_SIZE} (50)
schemes to balance throughput against memory usage.  A three-attempt
exponential back-off (delays: 2\,s, 4\,s, 8\,s) handles transient 5xx
errors and timeouts.

\subsubsection{Stage 3a — Benchmark Prices via yfinance
  (\texttt{fetch\_benchmark\_data})}

The 23 indices in \texttt{BENCHMARKS} (Table~\ref{tab:benchmarks}) are
downloaded via the \texttt{yfinance} library.  For each index the latest stored
date is retrieved and only the incremental tail is fetched using yfinance's
\texttt{start} parameter.  Close prices are stored under the index's display
name (e.g.\ \texttt{"Nifty 50"}).

\subsubsection{Stage 3b — NSE Archive Indices (\texttt{fetch\_nse\_index\_data})}

Seven indices — Nifty Smallcap 50/100/250, Nifty Microcap 250, Nifty Midcap
100, Nifty 500, and Nifty MidSmallcap 400 — are not available with meaningful
history on yfinance.  NSE publishes a bulk archive of all index closes in a
daily CSV:

\begin{center}
\texttt{archives.nseindia.com/content/indices/ind\_close\_all\_DDMMYYYY.csv}
\end{center}

Each file is approximately 16\,KB and contains every NSE index for that trading
day.  The downloader builds a reverse lookup from lowercase CSV column names to
display names, finds the earliest date not yet stored, generates the list of
weekday dates through today, downloads each file with a 50\,ms inter-request
delay, and performs bulk inserts.  Files for dates before an index existed in
the archive simply produce zero records for that index with no error.

\begin{center}
\begin{tabular}{@{}lll@{}}
\toprule
\textbf{Index} & \textbf{Source} & \textbf{History from} \\
\midrule
Nifty 50, Nifty Bank, Nifty IT, \ldots & yfinance & Varies (2000+) \\
Nifty Smallcap 50/100/250 & NSE archive CSV & Nov 2015 \\
Nifty Microcap 250 & NSE archive CSV & Jan 2016 \\
Nifty Midcap 100 & NSE archive CSV & Apr 2016 \\
Nifty 500 & NSE archive CSV & Jan 2013 \\
Nifty MidSmallcap 400 & NSE archive CSV & Nov 2015 \\
\bottomrule
\end{tabular}
\end{center}

\subsubsection{Incremental Update Logic}

All fetchers follow the same pattern: read the watermark (\texttt{MAX(date)})
for each series from the DB, then request only newer data.  First-time
population downloads the full history; subsequent runs fetch only the delta.
The \texttt{update.py} script orchestrates Stages 1--3 in sequence and can be
run on-demand or scheduled via cron.

% -----------------------------------------------------------
\subsection{Application Layer — Streamlit Pages}
% -----------------------------------------------------------

The UI is a single-file Streamlit application.  Navigation is via a sidebar
radio widget with five pages.

\subsubsection{Dashboard}

\textbf{Purpose}: Cross-fund returns heatmap with filtering.

\noindent\textbf{Sidebar controls}:
\begin{itemize}[noitemsep]
  \item Sector/style multiselect (from \texttt{fund\_tags})
  \item SEBI category multiselect
  \item Fund filter: Show all / Exclude funds / Pick funds
  \item Fund house, min history, and nav-count filters
\end{itemize}

\noindent\textbf{Main area}: A styled \texttt{pd.DataFrame} where each cell
shows a trailing return period (6M, 1Y, 2Y, 3Y, 5Y) for each fund.  Color
coding uses \texttt{\_signed\_gradient\_col}: the RdYlGn colormap is applied
symmetrically around zero so any positive return is green and any negative
return is red, per column.

\subsubsection{Fund Analysis}

\textbf{Purpose}: Deep dive into a single fund with a full analytics suite.

\noindent\textbf{Sidebar controls}: Fund picker, benchmark picker (with
\texttt{TAG\_BENCHMARK} auto-suggestion), date-range slider.

\noindent\textbf{Tabs}:
\begin{enumerate}[noitemsep]
  \item \textbf{NAV Chart} — line chart with benchmark overlay
  \item \textbf{Returns} — trailing return table, calendar year bar chart,
    SIP XIRR summary
  \item \textbf{Risk Metrics} — volatility, Sharpe, Sortino, Calmar, Martin,
    Ulcer, VaR, CVaR, downside deviation, max drawdown
  \item \textbf{Drawdown} — drawdown time series, drawdown statistics
  \item \textbf{Benchmark-Relative} — alpha, beta, $R^2$, tracking error,
    information ratio, capture ratios, batting average; rolling alpha and beta
    charts; rolling Sharpe/Sortino comparison
\end{enumerate}

All 45 metrics in the Risk tab are computed by a module-level
\texttt{@st.cache\_data}-decorated wrapper around \texttt{compute\_fund\_analytics}
so that changing a rolling-window widget does not re-run the full XIRR
root-finding computation.

\subsubsection{Fund Comparison}

\textbf{Purpose}: Side-by-side comparison of up to 10 user-selected funds.

\noindent\textbf{Tabs}:
\begin{enumerate}[noitemsep]
  \item \textbf{NAV Chart} — rebased to 100 from common start date
  \item \textbf{Returns} — returns heatmap across funds and periods
  \item \textbf{Rolling} — rolling return and rolling Sharpe over user-selected
    window
  \item \textbf{Full Table} — all trailing return columns sortable
  \item \textbf{Risk} — volatility, Sharpe, max drawdown, peer quartile rank
    for each fund
\end{enumerate}

\subsubsection{Sector \& Style}

\textbf{Purpose}: Peer group comparison within a tag (sector or style).

\noindent\textbf{Sidebar controls}:
\begin{itemize}[noitemsep]
  \item Tag picker (all tags from \texttt{fund\_tags})
  \item Benchmark picker with \texttt{TAG\_BENCHMARK} auto-selection
  \item Fund filter: Show all / Exclude funds / Pick funds
\end{itemize}

\noindent\textbf{Tabs}: Returns heatmap, risk metrics comparison, peer
distribution charts, tag management (assign/remove tags per fund).

\subsubsection{Data Management}

\textbf{Purpose}: Data pipeline controls and database statistics.

\noindent\textbf{Steps exposed in UI}:
\begin{enumerate}[noitemsep]
  \item Refresh scheme list from AMFI
  \item Download NAV history (full or incremental) with live progress
  \item Download yfinance benchmark data
  \item Download NSE archive indices (Smallcap, Microcap, Midcap 100, etc.)
  \item Re-run auto-tagging
  \item Daily update (all steps in sequence)
\end{enumerate}

DB statistics (total funds, NAV records, benchmarks loaded, last update
time) are displayed on the page.

% -----------------------------------------------------------
\subsection{Fund Tagging Rule Engine}
\label{sec:tagger}
% -----------------------------------------------------------

Every fund in the database receives one or more tags via \texttt{auto\_tag\_all\_funds()}.
The engine operates in two passes:

\paragraph{Pass 1 — Style from SEBI category.}
Non-sectoral equity funds are uniquely identified by their SEBI category string.
The \texttt{CATEGORY\_TAG} dict maps each category directly to a style tag:

\begin{center}
\begin{tabular}{@{}ll@{}}
\toprule
\textbf{SEBI Category} & \textbf{Tag} \\
\midrule
Equity Scheme - Large Cap Fund & Large Cap \\
Equity Scheme - Mid Cap Fund & Mid Cap \\
Equity Scheme - Small Cap Fund & Small Cap \\
Equity Scheme - Large \& Mid Cap Fund & Large \& Mid Cap \\
Equity Scheme - Multi Cap Fund & Multi Cap \\
Equity Scheme - Flexi Cap Fund & Flexi Cap \\
Equity Scheme - Focused Fund & Focused \\
Equity Scheme - Value Fund & Value \\
Equity Scheme - Contra Fund & Contra \\
Equity Scheme - ELSS & ELSS \\
\bottomrule
\end{tabular}
\end{center}

If a match is found, the function returns immediately — no name-based matching
is needed.

\paragraph{Pass 2 — Theme from fund name (sectoral/thematic only).}

For funds in the \texttt{"Equity Scheme - Sectoral/ Thematic"} SEBI category,
the engine compares the lowercased fund name against an ordered list of
\texttt{(tag, keywords)} rules (\texttt{THEME\_RULES}).  The \emph{first}
matching rule wins, so rule ordering is critical.  Rules for highly specific
or easily conflated themes (e.g.\ Business Cycle and Momentum) are placed
above the broader Quant \& Factor rule to prevent misclassification by AMC
name tokens like ``Quant''.

If no rule matches, the fund receives the tag \texttt{"Other"}.  Funds can
match at most one theme (the \texttt{break} after the first match enforces
this).

\begin{center}
\begin{tabular}{@{}lp{0.55\textwidth}@{}}
\toprule
\textbf{Tag} & \textbf{Sample keywords} \\
\midrule
Banking \& Financial Services & banking, financial services, bfsi \\
Technology & technology, digital india, information technology \\
Healthcare & pharma, healthcare, health and wellness \\
Infrastructure & infrastructure, build india, roads, railways \\
Consumption \& FMCG & consumption, fmcg, consumer durable \\
Manufacturing & manufacturing, make in india \\
Energy \& Resources & energy, natural resources, new energy \\
PSU & psu fund, public sector \\
Defence & defence, defense, strategic \\
ESG \& Ethical & esg, ethical fund, best-in-class \\
Business Cycle & business cycle \\
Momentum & momentum fund, active momentum \\
Quant \& Factor & quant fund, multi-factor, smart beta \\
International & global equity, japan equity, emerging market \\
\ldots & (24 tags total) \\
\bottomrule
\end{tabular}
\end{center}

% -----------------------------------------------------------
\subsection{Caching and Performance}
% -----------------------------------------------------------

\paragraph{Streamlit \texttt{@st.cache\_data}.}
All expensive computations are wrapped in module-level
\texttt{@st.cache\_data}-decorated functions.  The cache key is the function
arguments tuple; Streamlit serialises \texttt{pd.Series} inputs via their
hash.  The critical module-level wrappers are:

\begin{itemize}[noitemsep]
  \item \texttt{\_cached\_fund\_analytics(nav, bench)} — wraps
    \texttt{compute\_fund\_analytics}; avoids recomputing 45 metrics
    (including four XIRR root-finding operations) when the user changes a
    widget that does not affect the fund or date range.
  \item \texttt{\_cached\_rolling\_sip(scheme\_code, yrs)} — wraps
    \texttt{rolling\_sip\_xirr}; the rolling 3-year SIP XIRR is an $O(N)$
    scan of XIRR calls and can take 2--5\,s for funds with 10+ year histories.
  \item \texttt{load\_nav\_as\_series(scheme\_code)} — database read cached by
    scheme code; invalidated on \texttt{st.cache\_data.clear()} after any
    data download.
\end{itemize}

\textbf{Important}: cached functions must be defined at module scope, not
inside \texttt{with} blocks or nested function bodies.  Locally-defined
decorated functions receive a new identity on each Streamlit rerun, making
the cache ineffective.

\paragraph{SQLite WAL mode.}
The database is opened with \texttt{PRAGMA journal\_mode=WAL} and
\texttt{PRAGMA synchronous=NORMAL} on every connection.  WAL mode allows
concurrent reads while a write is in progress, which prevents the Streamlit UI
from blocking during background data downloads.

\paragraph{Incremental fetching.}
Neither NAV history nor benchmark prices are ever re-downloaded in full after
the initial population.  Every fetch call reads \texttt{MAX(date)} from the DB
and requests only the tail.  For daily updates this means each scheme requires
a single API call returning typically 1--5 records.

% -----------------------------------------------------------
\subsection{Configuration Reference}
% -----------------------------------------------------------

All tuneable constants live in \texttt{config.py}.
Table~\ref{tab:config} documents every constant.

\begin{longtable}{@{}p{0.32\textwidth}p{0.15\textwidth}p{0.43\textwidth}@{}}
\caption{config.py constants}\label{tab:config}\\
\toprule
\textbf{Constant} & \textbf{Default} & \textbf{Purpose} \\
\midrule
\endfirsthead
\multicolumn{3}{c}{\textit{(continued)}}\\
\toprule
\textbf{Constant} & \textbf{Default} & \textbf{Purpose} \\
\midrule
\endhead
\texttt{DB\_PATH} & \texttt{./mf\_data.db} & SQLite file path \\[3pt]
\texttt{AMFI\_NAV\_URL} & amfiindia.com/\ldots & AMFI scheme list endpoint \\[3pt]
\texttt{MFAPI\_BASE\_URL} & api.mfapi.in/mf & NAV history base URL \\[3pt]
\texttt{MFAPI\_DELAY\_SECONDS} & 0.5 & Inter-request delay for mfapi.in \\[3pt]
\texttt{MFAPI\_BATCH\_SIZE} & 50 & DB commit frequency during NAV download \\[3pt]
\texttt{RISK\_FREE\_RATE} & 0.065 & Annualised risk-free rate (India 91-day T-bill) \\[3pt]
\texttt{TRADING\_DAYS\_PER\_YEAR} & 252 & Used for annualisation of daily statistics \\[3pt]
\texttt{BENCHMARKS} & 23 entries & Dict of display name $\to$ yfinance ticker \\[3pt]
\texttt{NSE\_DIRECT\_INDICES} & 7 entries & Dict of display name $\to$ NSE CSV column name \\[3pt]
\texttt{NSE\_ARCHIVE\_URL} & archives.nseindia.com/\ldots & URL template for NSE bulk CSV \\[3pt]
\texttt{TAG\_BENCHMARK} & $\sim$30 entries & Preferred benchmark per sector/style tag \\[3pt]
\texttt{EQUITY\_CATEGORY\_KEYWORDS} & 2 entries & Substrings to detect equity/ELSS funds \\[3pt]
\texttt{ROLLING\_WINDOWS} & 5 entries & Name $\to$ trading-day count (6M=126, \ldots, 5Y=1260) \\[3pt]
\texttt{VAR\_CONFIDENCE\_LEVELS} & [0.95, 0.99] & Confidence levels for VaR / CVaR \\
\bottomrule
\end{longtable}

% -----------------------------------------------------------
\subsection{Benchmark Coverage}
% -----------------------------------------------------------

\begin{longtable}{@{}p{0.38\textwidth}p{0.28\textwidth}p{0.24\textwidth}@{}}
\caption{Benchmark indices available in the platform}\label{tab:benchmarks}\\
\toprule
\textbf{Index} & \textbf{Source} & \textbf{Category} \\
\midrule
\endfirsthead
\multicolumn{3}{c}{\textit{(continued)}}\\
\toprule
\textbf{Index} & \textbf{Source} & \textbf{Category} \\
\midrule
\endhead
Nifty 50 & yfinance & Broad market \\
Nifty 100 & yfinance & Broad market \\
BSE Sensex & yfinance & Broad market \\
Nifty Next 50 & yfinance & Broad market \\
Nifty Midcap 50 & yfinance & Cap segment \\
Nifty Midcap 150 & yfinance & Cap segment \\
Nifty Smallcap 50 & NSE archive & Cap segment \\
Nifty Smallcap 100 & NSE archive & Cap segment \\
Nifty Smallcap 250 & NSE archive & Cap segment \\
Nifty Microcap 250 & NSE archive & Cap segment \\
Nifty Midcap 100 & NSE archive & Cap segment \\
Nifty MidSmallcap 400 & NSE archive & Cap segment \\
Nifty 500 & NSE archive & Broad market \\
Nifty Bank & yfinance & Sector \\
Nifty PSU Bank & yfinance & Sector \\
Nifty Financial Svc & yfinance & Sector \\
Nifty IT & yfinance & Sector \\
Nifty FMCG & yfinance & Sector \\
Nifty Consumption & yfinance & Sector \\
Nifty Pharma & yfinance & Sector \\
Nifty Infrastructure & yfinance & Sector \\
Nifty Energy & yfinance & Sector \\
Nifty Metal & yfinance & Sector \\
Nifty Media & yfinance & Sector \\
Nifty Commodities & yfinance & Sector \\
Nifty Services & yfinance & Sector \\
Nifty PSE & yfinance & Sector \\
Nifty Auto & yfinance & Sector \\
Nifty Realty & yfinance & Sector \\
Nifty MNC & yfinance & Sector \\
\bottomrule
\end{longtable}

% -----------------------------------------------------------
\subsection{Deployment and Setup}
% -----------------------------------------------------------

\subsubsection{Requirements}

\begin{itemize}[noitemsep]
  \item Python 3.11+ (tested on 3.14)
  \item Packages: \texttt{streamlit}, \texttt{pandas}, \texttt{numpy},
    \texttt{scipy}, \texttt{plotly}, \texttt{yfinance}, \texttt{requests},
    \texttt{matplotlib}
  \item Optional: \texttt{schedule} (for \texttt{update.py --schedule} mode)
\end{itemize}

\subsubsection{First-time setup}

\begin{enumerate}[noitemsep]
  \item \texttt{pip install -r requirements.txt}
  \item \texttt{python update.py} \quad — populates scheme list and downloads
    NAV history for all equity Direct Growth funds ($\approx$3--6\,h for the
    initial full download of $\sim$1{,}700 schemes)
  \item \texttt{streamlit run app.py} \quad — launches the web UI on
    \texttt{localhost:8501}
  \item In the Data Management page: download yfinance benchmarks, then NSE
    archive indices, then run auto-tagging
\end{enumerate}

\subsubsection{Daily update}

The simplest approach is a \texttt{cron} job:

\begin{lstlisting}[style=pycode,language=bash,caption={Cron entry for weekday updates at 22:00 IST}]
0 22 * * 1-5  cd /path/to/mf_analytics && python update.py >> update.log 2>&1
\end{lstlisting}

Alternatively, the built-in scheduler mode keeps a persistent process running:

\begin{lstlisting}[style=pycode,language=bash,caption={Persistent scheduler mode}]
python update.py --schedule --time 22:00
\end{lstlisting}

The Streamlit UI also exposes a ``Daily Update'' button in the Data Management
page for on-demand manual updates with a live progress indicator.

\subsubsection{Data flow summary}

\begin{center}
\begin{tabular}{@{}lll@{}}
\toprule
\textbf{Action} & \textbf{Duration (first run)} & \textbf{Duration (daily)} \\
\midrule
Scheme list refresh & $<$5\,s & $<$5\,s \\
NAV history download & 3--6\,h ($\sim$1700 schemes) & 5--15\,min \\
yfinance benchmarks & $<$2\,min & $<$30\,s \\
NSE archive indices & 15--30\,min (3453 files) & $<$5\,min \\
Auto-tagging & $<$1\,s & $<$1\,s \\
\bottomrule
\end{tabular}
\end{center}

% ============================================================
\section{References}
% ============================================================

\begin{enumerate}[label={[\arabic*]}]

\item Sharpe, W. F. (1966). ``Mutual Fund Performance.''
  \textit{Journal of Business}, 39(1), 119--138.

\item Jensen, M. C. (1968). ``The Performance of Mutual Funds in the
  Period 1945--1964.''
  \textit{Journal of Finance}, 23(2), 389--416.

\item Sortino, F. A., \& Van der Meer, R. (1991). ``Downside Risk.''
  \textit{Journal of Portfolio Management}, 17(4), 27--31.

\item Young, T. W. (1991). ``Calmar Ratio: A Smoother Tool.''
  \textit{Futures}, 20(1), 40.

\item Martin, P., \& McCann, B. (1989).
  \textit{The Investor's Guide to Fidelity Funds}.
  John Wiley \& Sons. (Ulcer Index originally defined here.)

\item Treynor, J. L. (1965). ``How to Rate Management of Investment
  Funds.'' \textit{Harvard Business Review}, 43(1), 63--75.

\item Artzner, P., Delbaen, F., Eber, J.-M., \& Heath, D. (1999).
  ``Coherent Measures of Risk.''
  \textit{Mathematical Finance}, 9(3), 203--228.
  (Formal basis for CVaR / Expected Shortfall.)

\item Grinold, R. C., \& Kahn, R. N. (2000).
  \textit{Active Portfolio Management} (2nd ed.).
  McGraw-Hill. (Information Ratio, tracking error.)

\item Goodwin, T. H. (1998). ``The Information Ratio.''
  \textit{Financial Analysts Journal}, 54(4), 34--43.

\item AMFI India. (2024). \textit{NAVAll.txt specification}.
  \url{https://www.amfiindia.com/spages/NAVAll.txt}

\item SEBI. (2017). \textit{Circular on Categorization and Rationalization
  of Mutual Fund Schemes}. SEBI/HO/IMD/DF3/CIR/P/2017/114.

\end{enumerate}

\vfill
\begin{center}
\rule{0.5\textwidth}{0.4pt}\\[0.5em]
{\small\textcolor{darkgray}{
  MF Analytics Platform \textbullet{}
  All calculations implemented in \texttt{analytics.py} \textbullet{}
  Risk-free rate 6.5\% p.a. \textbullet{}
  252 trading days per year
}}
\end{center}

\end{document}
